% Chapter 9, typeset from lectures/L09/README.md.

\chapter{Practical exam 1: sequential networks}
\label{c9:ch}

{\large\itshape Three hours, individual, at the computer. Covers the whole of part~\ref{part:vhdl},
with the emphasis on sequential networks.\par}

\begin{chapterbox}{In this chapter}
\begin{itemize}
  \item The exam's format: independent tasks, each with a self-checking testbench handed out.
  \item What you are allowed to use, and what you are not.
  \item How to prepare, and what the best preparation there is.
  \item What the exam actually measures.
\end{itemize}
\end{chapterbox}

This is the course's first examination. It contains no new theory; the chapter says what the exam
measures, what shape the tasks take and how to prepare, and it is worth reading in advance rather
than the day before. Marks and grade boundaries are in appendix~\ref{app:project}.

The exam covers \chapref{c1:ch} to \ref{c8:ch} in full, with the emphasis on sequential networks:
registers, counters, shift registers, timers, synchronisers and simple Moore machines.

% ------------------------------------------------------------------------------------------
\section{Format}
\label{c9:sec:format}

The exam consists of a number of independent tasks. Every task gives you a \textbf{self-checking
testbench, handed out}, and a description of the module it drives: port order, types and behaviour.
Your job is to write the module so that the testbench passes.

\begin{itemize}
  \item The tasks are independent of each other. A module you cannot get working stops no other.
  \item The tasks are ordered roughly by increasing difficulty, but read all of them before you
    start.
  \item A module that is not entirely finished is marked on what is there. Partially correct logic
    earns partial marks, so never hand in an empty file because it did not get finished.
  \item The testbench is the answer key during the exam, not after it. \code{all checks passed!}
    means that task is done; the marking reads the code as well.
\end{itemize}

A practice paper with tasks of the same kind and the same scope is handed out in advance.

% ------------------------------------------------------------------------------------------
\section{Permitted aids}
\label{c9:sec:aids}

\textbf{Allowed:}
\begin{itemize}
  \item All the course material: this book, the lectures' READMEs, every appendix, and all the VHDL
    you wrote yourself during part~\ref{part:vhdl}, including your solutions to the exercises.
  \item Your own computer, with GHDL. See \secref{c2:sec:testbenches} if you need reminding how a
    testbench is run.
  \item CircuitVerse, if you want to draw a circuit before you write it. That is often faster, not
    slower.
\end{itemize}

\textbf{Not allowed:} communication with others in any form; language models and other tools that
generate or explain VHDL for you; and searching the internet for ready-made solutions. The line is
not drawn at who presses the keys but at whose understanding is being tested. Looking up GHDL's or
VHDL's documentation is allowed; searching for a finished module is not.

The full and authoritative list of permitted aids is in appendix~\ref{app:project}; if anything here
says otherwise, that one applies.

% ------------------------------------------------------------------------------------------
\section{Before the exam}
\label{c9:sec:prep}

\begin{itemize}
  \item Work through the practice paper, under time pressure and without looking at the solutions
    first.
  \item Exercise~\ref{c8:ex:rx} and \ref{c8:ex:tx}, the book's two capstones, are the best
    preparation there is: they force you to design a state diagram yourself and compose earlier
    modules around it, which is exactly what the exam measures.
  \item Make sure GHDL works on your machine \emph{before} you arrive. A broken toolchain is no
    reason for extra time.
\end{itemize}

% ------------------------------------------------------------------------------------------
\section{What the exam measures}
\label{c9:sec:measures}

\begin{itemize}
  \item That you can write the synchronous process template correctly, without looking it up.
  \item That you can build a counter, a timer and a shift register and know what separates them.
  \item That you synchronise every asynchronous input, without being reminded.
  \item That you can design a small Moore machine from a specification in words, and write it as an
    enumerated type and a \code{case}.
  \item That you can read a testbench's error output and find your way back to the line that caused
    it.
\end{itemize}

% ------------------------------------------------------------------------------------------
\section{After the exam}
\label{c9:sec:after}

The next chapter starts the \textbf{group project}, which runs from \chapref{c10:ch} to
\ref{c19:ch} and accounts for half the course's marks. Read the project specification in
appendix~\ref{app:project} before \chapref{c10:ch}; the groups are put together there.

\Chapref{c10:ch} takes up the project start: the CAN bus, the frame and arbitration, and the
project's first VHDL, the shared package \code{can_def}.
